% paper5_discussion.tex -- Discussion and Conclusion of Paper 5, input by
% papers/paper5_amendments.tex.
\section{Discussion}

Most of the carbon added in the amendments was lost from the soil within five
years, and the fraction retained was small at every site. This agrees with
the view that the persistence of organic matter is controlled less by the
chemistry of the material added than by the protection it receives once in the
soil \parencite{vonlutzow2006}. The higher retention of compost than of fresh
residues is consistent with that view: compost has already lost its most
readily decomposed compounds during composting, so a larger share of what
remains is available for association with mineral surfaces.

Retention increased with clay content across the three sites. Clay and fine
silt provide the surfaces on which organic matter is stabilised and the
microaggregates within which it is occluded, so soils with more clay can
protect a larger share of each addition \parencite{six2002}. The relationship
has a practical implication: the same amendment will build soil carbon more
efficiently on a clay soil than on a sandy one, and claims for carbon
sequestration from amendments should account for soil texture. With only three
sites, however, texture is confounded with climate in this trial, and the
semi-arid site is both the sandiest and the driest.

Amendments add carbon that was fixed elsewhere, so their contribution to
removing carbon dioxide from the atmosphere is smaller than the gain in soil
carbon suggests; part of the benefit is a transfer from one place to another.
Their value for soil health\index{soil health} and crop nutrition remains, and they form a core
component of integrated soil fertility management \parencite{vanlauwe2010}.

\section{Conclusion}

Averaged over three amendments, the fraction of added carbon retained after
five years was \res{p5-ret-S2}, \res{p5-ret-S1} and \res{p5-ret-S3} at the
sites with \res{site-clay-S2}, \res{site-clay-S1} and \res{site-clay-S3} per cent clay. Compost was retained best and crop
residues least at every site. Texture should be considered whenever the carbon
benefit of an amendment is estimated. The data are synthetic and illustrate the
calculation only.
